% Capitulo7_Conclusiones.tex
\chapter{Conclusiones y Trabajo Futuro}
\label{cap:conclusiones}

\section{Limitaciones del Trabajo}
A pesar de la validez del marco metodológico desarrollado y la consistencia de los resultados obtenidos, es necesario explicitar las siguientes limitaciones asumidas en la presente investigación:

\begin{enumerate}
    \item \textbf{Ausencia de Penalización en el Vector de Control Dinámico de Vuelo:} 
    Las restricciones físicas de no transitabilidad (masas de agua profunda y edificaciones urbanas) han sido implementadas como una \textit{restricción dura de probabilidad inicial} ($P_0(x,y) = 0.0$), garantizando que los algoritmos de búsqueda guiados por información no concentren recursos en dichas zonas. Sin embargo, no se ha incorporado un término de penalización o barrera física en las acciones del vector de control dinámico del UAV, lo que permite que el dron pueda sobrevolar físicamente dichas zonas no transitables en su trayectoria de desplazamiento hacia otras celdas de alto interés.

    \item \textbf{Asunción de Objetivos Estáticos:} 
    Todas las simulaciones y escenarios del benchmark consideran que la persona desaparecida permanece en una ubicación fija e inalterable desde el inicio de la operación de búsqueda hasta su hallazgo o el agotamiento del presupuesto temporal.
\end{enumerate}

\section{Líneas de Trabajo Futuro}
Como extensión natural del trabajo presentado, se proponen las siguientes líneas de desarrollo futuro:

\begin{itemize}
    \item \textbf{Modelado de Víctimas Dinámicas (Cadenas de Markov):} 
    Incorporación de modelos dinámicos de movimiento para personas desorientadas en función del perfil de comportamiento, utilizando matrices de transición de Markov orientadas por la pendiente del terreno y la conectividad de los senderos.
    \item \textbf{Control de Vuelo Guiado por Restricciones Geométricas:} 
    Implementación de campos potenciales de repulsión o algoritmos de evitación de zonas prohibidas en la capa de control cinemático del dron para impedir el sobrevuelo físico de masas de agua.
\end{itemize}
